\documentclass[11pt,a4,oneside]{article}

\usepackage{graphicx} % Required for inserting images

\usepackage{geometry}

\usepackage{longtable}

\usepackage{multirow}

\usepackage{caption}

\usepackage{subcaption}

\usepackage[table,xcdraw]{xcolor}

\usepackage{bibentry}

\usepackage[skip=10pt plus1pt, indent=40pt]{parskip}

\usepackage[portuguese]{babel}

\geometry{
a4paper,
left=3cm,
top=3cm,
right=2cm,
bottom=2cm
}

\usepackage{fancyhdr}

\usepackage{biblatex} %Imports biblatex package

\addbibresource{references.bib} %Import the bibliography file

\linespread{1.2}

\title{\textit{Checklist} para verificação de aderência do tema de Trabalho de COnclusão de Curso I ao curso de Ciência da Computação.}

\date{Agosto 2023}

\begin{document}

\pagestyle{fancy}

%\fancyhead[C]{In the centre of the header on all pages}

\pagestyle{fancy}

\fancyhead{}

\fancyhead[RO,LE]{\textbf{Checklist - TCC I}}

\maketitle

\begin{itemize}

\item \textbf{Modalidade:} Acadêmico - Projeto - Grupo

\item \textbf{Tema:} Correlação entre o formato de uma prova e o desempenho do aluno

\item \textbf{Alunos:}

\begin{itemize}

\item Beatriz Manaia Lourenço Berto (RA: 22.125.060-8)

\item Luana Bortko Rodrigues (RA:24.123.006-9)

\item Nuno Martins Guilhermino da Silva (RA:22.126.099-5)

\end{itemize}

\item \textbf{Orientador:} Prof.Charles Henrique Porto Ferreira

\item \textbf{Resumo:}

O formato das questões do Exame Nacional do Ensino Médio (ENEM) possui características que influenciam o desempenho dos alunos. Aspectos como extensão, grau de formalismo da linguagem e nível de dificuldade dos itens — à luz da Teoria de Resposta ao Item (TRI) — exercem impacto direto sobre o resultado obtido. Assim, as notas obtidas pelos alunos não refletem o domínio do conteúdo, mas como estudantes lidam com as características estruturais da prova. Este trabalho propõe analisar os formatos das provas do ENEM ao longo dos anos, a fim de identificar quais fatores apresentam maior ou menor influência sobre a nota final dos estudantes. Faremos isso através de um programa que contará o número de palavras contidas na prova e nas questões para que possamos identificar padrões, também consultaremos as informações sobre o  TRI divulgadas pelo INEP. 


\end{itemize}

\section*{Lista de disciplinas aderentes} analise de dados, ciencia de dados, estatistica,

\begin{itemize}

\item \textbf{CC1612 Fundamentos de Algoritmos: }Utilizaremos conceitos de lógica e programação básica, como laços de repetição, para o desenvolvimento desse trabalho.

\item \textbf{CC2632 Desenvolvimento de Algoritmos: } :D

\item \textbf{CC3642 Orientação A Objetos}: Usaremos conceitos os pilares da programação orientada a objetos, como polimorfismo e herança.

\item \textbf{CC4652 Estrutura de dados}: O uso de diferentes estruturas aprendidas na disciplina auxiliará na manipulação e armazenamento de dados encontrados.

\item \textbf{CA5411 Probabilidade e Estatística}: Fórmulas aprendidas durante esta disciplina são relevantes para a análise dos dados do projeto.

\item \textbf{CC5561 Análise de Complexidade de Algoritmos}: As técnicas para otimização e análise de programas aprendidas resultarão em um código melhor estruturado.

\item \textbf{CC7711 Inteligência Artificial e Robótica}: Serão utilizadas técnicas de inteligência artificial para encontrar padrões nas questões.

\end{itemize}

\end{document}\documentclass[11pt,a4,oneside]{article}

\usepackage{graphicx} % Required for inserting images

\usepackage{geometry}

\usepackage{longtable}

\usepackage{multirow}

\usepackage{caption}

\usepackage{subcaption}

\usepackage[table,xcdraw]{xcolor}

\usepackage{bibentry}

\usepackage[skip=10pt plus1pt, indent=40pt]{parskip}

\usepackage[portuguese]{babel}

\geometry{

a4paper,

left=3cm,

top=3cm,

right=2cm,

bottom=2cm

}

\usepackage{fancyhdr}

\usepackage{biblatex} %Imports biblatex package

\addbibresource{references.bib} %Import the bibliography file

\linespread{1.2}

\title{\textit{Checklist} para verificação de aderência do tema de Trabalho de COnclusão de Curso I ao curso de Ciência da Computação.}

\date{Agosto 2023}

\begin{document}

\pagestyle{fancy}

%\fancyhead[C]{In the centre of the header on all pages}

\pagestyle{fancy}

\fancyhead{}

\fancyhead[RO,LE]{\textbf{Checklist - TCC I}}

\maketitle

\begin{itemize}

\item \textbf{Modalidade:} Acadêmico - Projeto - Grupo

\item \textbf{Tema:} Correlação entre o formato de uma prova e o desempenho do aluno

\item \textbf{Alunos:}

\begin{itemize}

\item Beatriz Manaia Lourenço Berto (RA: 22.125.060-8)

\item Luana Bortko Rodrigues (RA:24.123.006-9)

\item Nuno Martins Guilhermino da Silva (RA:22.126.099-5)

\end{itemize}

\item \textbf{Orientador:} Prof.Charles Henrique Porto Ferreira

\item \textbf{Resumo:}

O formato das questões do Exame Nacional do Ensino Médio (ENEM) possui características que influenciam o desempenho dos alunos. Aspectos como extensão, grau de formalismo da linguagem e nível de dificuldade dos itens — à luz da Teoria de Resposta ao Item (TRI) — exercem impacto direto sobre o resultado obtido. Assim, as notas obtidas pelos alunos não refletem o domínio do conteúdo, mas como estudantes lidam com as características estruturais da prova. Este trabalho propõe analisar os formatos das provas do ENEM ao longo dos anos, a fim de identificar quais fatores apresentam maior ou menor influência sobre a nota final dos estudantes. ADICIONAR COMO VAMOS FAZER!!!!!!!!!!!!!

\end{itemize}

\section*{Lista de disciplinas aderentes} analise de dados, ciencia de dados, estatistica,

\begin{itemize}

\item \textbf{CC1612 Fundamentos de Algoritmos: }Utilizaremos conceitos de lógica e programação básica, como laços de repetição, para o desenvolvimento desse trabalho.

\item \textbf{CC2632 Desenvolvimento de Algoritmos: } :D

\item \textbf{CC3642 Orientação A Objetos}: Usaremos conceitos os pilares da programação orientada a objetos, como polimorfismo e herança.

\item \textbf{CC4652 Estrutura de dados}: O uso de diferentes estruturas aprendidas na disciplina auxiliará na manipulação e armazenamento de dados encontrados.

\item \textbf{CA5411 Probabilidade e Estatística}: Fórmulas aprendidas durante esta disciplina são relevantes para a análise dos dados do projeto.

\item \textbf{CC5561 Análise de Complexidade de Algoritmos}: As técnicas para otimização e análise de programas aprendidas resultarão em um código melhor estruturado.

\item \textbf{CC7711 Inteligência Artificial e Robótica}: Serão utilizadas técnicas de inteligência artificial para encontrar padrões nas questões.

\end{itemize}

\end{document}